\section{Design Discussion: ``Method'' vs.\ ``Algorithm'' Granularity}
\label{sec:appendix-method}

A referee raised a subtle but important point: what the MLIP community
calls an ``algorithm'' is not, strictly speaking, an algorithm in the
ML-Schema sense. This appendix documents the granularity mismatch, the
terminology choice, and the design we have adopted --- an
$\MLIPMethodC$ class that is disjoint from $\mlsAlgorithmC$ and is
decomposed into its three conceptual components, with the training
algorithm contributed as the sole alignment point with ML-Schema.

\subsection{The mismatch}

In ML-Schema, $\mlsAlgorithmC$ is defined as ``a computational
procedure''---typically a training algorithm such as stochastic gradient
descent, Adam, or L-BFGS. The procedure consumes data and hyperparameters
and produces a fitted $\mlsModelC$. This is the classical ML view:
\emph{model family} and \emph{training procedure} are conceptually
distinct, and $\mlsAlgorithmC$ names only the latter.

In the MLIP community, a named ``algorithm'' (MACE, MTP, ACE, NequIP)
typically bundles \emph{four} different things:
%
\begin{enumerate}
\item a \textbf{descriptor} of the local atomic environment (moment
  tensors, SOAP, symmetry functions, equivariant messages);
\item a \textbf{functional form} that maps descriptor features to the
  local energy (linear expansion in basis functions, neural network,
  kernel regression);
\item a \textbf{loss function} for training (typically a weighted
  combination of energy error, force error, and stress error); and
\item a \textbf{training procedure} in the ML-Schema sense (an
  optimizer, sampling scheme, and stopping criterion).
\end{enumerate}
%
Only item~(4) aligns directly with $\mlsAlgorithmC$. Items (1)--(3) are
properties of the \emph{model family} rather than of any training
procedure, and the bundle as a whole is what the community refers to
by a single name.

\subsection{Analogy: knowledge graph embeddings}

The same mismatch arises in other machine-learning sub-communities.
Consider TransE~\cite{bordes2013transe}, a standard knowledge graph
embedding method. ``TransE'' as used in the KGE literature bundles:
%
\begin{itemize}
\item an entity representation (points in $\mathbb{R}^d$);
\item a relation representation (translation vectors in $\mathbb{R}^d$);
\item a scoring function $f(h,r,t) = \|h + r - t\|$; and
\item a training procedure (stochastic gradient descent with
  margin-based ranking loss and negative sampling).
\end{itemize}
%
Asked ``what is TransE in ML-Schema?'', a strict reading splits it
across $\mlsAlgorithmC$ (the SGD procedure) and $\mlsModelC$ (the fitted
embeddings). ML-Schema has no single class for the bundle that the
community actually names and identifies. The same structural issue
surfaces in graph neural networks (GCN, GAT, GIN), generative models
(VAE, GAN), and many other ML sub-communities where a named
``algorithm'' is a bundle of model-family choices plus a training
recipe.

\subsection{Terminology: ``Algorithm'', ``Method'', or ``Formalism''?}

Given that MTP and its siblings are \emph{not} algorithms in
ML-Schema's sense, it is worth asking which of three candidate words is
most accurate and most widely used.

\paragraph{Algorithm.}
Narrowly, an algorithm is a computational procedure --- a sequence of
steps. MTP is not a sequence of steps; it is a mathematical framework
that defines what energy function is being fitted plus the machinery to
fit it. ``The MTP algorithm'' is used casually in the literature but
is semantically imprecise.

\paragraph{Method.}
Generic: ``a named approach or recipe''. MTP is clearly a method: a
recipe for constructing a potential from a particular descriptor
(moment tensors), a particular functional form (linear expansion in
basis functions), a particular loss (weighted energy/force/stress),
and a particular training algorithm. Semantically clean; carries no
specific structural commitment beyond ``named bundle''.

\paragraph{Formalism.}
Emphasises the mathematical framework. In physics and chemistry,
``formalism'' is the standard word for ``a mathematical structure that
defines what a class of objects looks like and how they behave'' --- e.g.,
``the Lagrangian formalism'', ``the density functional theory
formalism'', ``the Standard Model formalism''. MTP fits this usage
well: the moment tensor expansion defines what MTP potentials are,
before any specific one is fitted. This is the tightest of the three
for MTP-the-mathematical-framework but carries a physics-flavoured
connotation.

\paragraph{Literature usage across MLIP papers.}
``Method'' dominates in practice (Novikov et al.\ 2021, Batatia et al.\
2022, Bart\'ok et al.\ 2010). ``Formalism'' and ``framework'' appear
when authors want to emphasise the mathematical structure (Drautz 2019
for ACE, and occasionally in the MTP literature). ``Algorithm'' is used
casually but rarely in formal technical contexts.

\paragraph{Cross-community applicability.}
For an upper ontology targeting multiple ML sub-communities, ``Method''
is the most portable label: it reads naturally for MLIPs, KGE, GNNs,
generative models, and classical ML alike. ``Formalism'' is accurate
for physics-adjacent fields but heavy in KGE (``the TransE formalism''
is correct but uncommon) and uncommon in GNN/CV/NLP. ``Algorithm'' is
semantically wrong everywhere.

\paragraph{Decision.}
We adopt \emph{Method} as the class name at the domain layer
($\MLIPMethodC$). Author freedom to use ``formalism'' or ``framework''
in prose when the mathematical-structure emphasis is appropriate is
retained; the class name does not prescribe author vocabulary.

\subsection{The design adopted}
\label{sec:appendix-method-design}

Rather than align $\MLIPMethodC$ with $\mlsAlgorithmC$ by subsumption
--- which would conflate the named bundle with the training procedure
it uses --- we declare them \emph{disjoint} and decompose the bundle
into its three non-training components, with the training algorithm
contributed as the sole alignment point with ML-Schema:
%
\begin{align*}
  \MLIPMethodC \sqcap \mlsAlgorithmC &\sqsub \bot \\
  \MLIPMethodC &\sqsub \existsR{hasFunctionalForm}{FunctionalForm} \\
  \MLIPMethodC &\sqsub \existsR{hasLossFunction}{LossFunction} \\
  \MLIPMethodC &\sqsub \existsR{hasTrainingAlgorithm}{\mlsAlgorithmC}
\end{align*}
%
Here $\mlsAlgorithmC$ is used in its narrow ML-Schema sense (the
optimiser/training procedure only). The bundle $\MLIPMethodC$ is not a
subclass of $\mlsAlgorithmC$; instead, one of its components \emph{is}
an $\mlsAlgorithmC$. The atomic-environment descriptor is modelled
separately via $\hasDescriptorR$, alongside the method's hyperparameters,
implementations, and supported simulation types.

Concrete instances then look like:
%
\begin{lstlisting}[language=turtle]
ex:MACE a mlips:MLIPMethod ;
    rdfs:label "MACE" ;
    mlips:hasDescriptor mlips:EquivariantMessagePassingDescriptor ;
    mlips:hasFunctionalForm ex:MACE-NeuralNet ;
    mlips:hasLossFunction ex:WeightedEnergyForceStressLoss ;
    mlips:hasTrainingAlgorithm ex:Adam ;
    mlips:hasHyperparameter ex:hp-rcut, ex:hp-num-layers ;
    mlips:supportsSimulation ex:sim-md .

ex:Adam a mls:Algorithm .  # narrow ML-Schema sense
\end{lstlisting}

For the same reason of disjointness, $\MLIPRunC$ is a subclass of
$\provActivityC$ only, not of $\mlsRunC$: $\mlsRunC.\mlsExecutesR$
ranges over $\mlsAlgorithmC$, which is disjoint from $\MLIPMethodC$. A
dedicated role $\appliesMethodR$ replaces $\mlsExecutesR$ for the MLIP
case. An optional $\hasTrainingRunR$ links an $\MLIPRunC$ to an
$\mlsRunC$ when the training algorithm is itself modelled in
ML-Schema.

\paragraph{Benefits.}
\begin{itemize}
\item The alignment with ML-Schema is tight and semantically correct:
  only the training procedure is an $\mlsAlgorithmC$.
\item Each component is independently queryable and comparable across
  methods: ``which methods share the same descriptor?'', ``which
  methods use a margin-based ranking loss?''
\item The pattern generalises: TransE, GCN, VAE, and other
  community-named bundles can be modelled analogously in sibling
  domain ontologies.
\end{itemize}

\paragraph{Costs.}
\begin{itemize}
\item MLIP researchers commonly leave the training procedure
  unspecified (it is often a detail not worth naming), so in practice
  many $\MLIPMethodC$ instances carry only a placeholder
  $\hasTrainingAlgorithmR$ link. The value of making it explicit is
  mostly conceptual, but is what enables the alignment with
  ML-Schema to be both tight and narrow.
\item The functional-form and loss-function components are often also
  left implicit in the literature; the ontology makes them first-class
  but does not require extensive axiomatisation of their internal
  structure.
\end{itemize}
